\chapter{مسئلۀ لگاریتم گسسته}
\label{app:dlp}
\section{تعریف و نقش در رمزنگاری}
در گروه چرخه‌ای متناهی $G=\langle g\rangle$ با مرتبۀ $q$، اگر $h=g^x$ معلوم باشد، مسئلۀ لگاریتم گسسته یافتن $x\in\Z_q$ است:
\begin{equation}
\operatorname{DLP}(g,h):\qquad g^x=h.
\end{equation}
توان‌رسانی با الگوریتم مربع‌کردن و ضرب در تعداد لگاریتمیِ عملیات گروه انجام می‌شود؛ یافتن وارون این نگاشت، بسته به ساختار و اندازۀ گروه، می‌تواند دشوار باشد. دشواری صرفاً به تعداد بیت وابسته نیست. مرتبۀ گروه، اندازۀ بزرگ‌ترین عامل اول مرتبه و نمایش گروه بر الگوریتم‌های حمله اثر می‌گذارند.

جست‌وجوی ساده تا $q$ مقدار، هزینۀ $O(q)$ دارد. روش \en{baby-step/giant-step} با زمان و حافظۀ $O(\sqrt q)$ و روش‌های تصادفی مانند \en{Pollard rho} با حافظۀ کمتر، مبادلۀ دیگری فراهم می‌کنند. اگر مرتبه به عامل‌های کوچک تجزیه شود، روش \en{Pohlig\lr{--}Hellman} مسئله را به زیرمسئله‌ها کاهش می‌دهد. بنابراین انتخاب گروه با مرتبۀ نامناسب می‌تواند امنیت را کاهش دهد، حتی اگر پیمانه در ظاهر بزرگ باشد.

\section{روش گام کوچک\lr{--}گام بزرگ}
با $m=\lceil\sqrt q\rceil$، نمای مجهول را به صورت $x=im+j$ می‌نویسیم. جدول گام کوچک شامل $g^j$ برای $0\le j<m$ ساخته می‌شود. سپس برای $i$های متوالی، مقدار $h(g^{-m})^i$ با جدول مقایسه می‌شود. هر برخورد معتبر می‌دهد:
\begin{equation}
h g^{-im}=g^j\quad\Longrightarrow\quad x=im+j\pmod q.
\end{equation}
بهتر است جواب بازیابی‌شده دوباره در رابطۀ $g^x=h$ بررسی شود. عملیات وارون نیازمند معتبر بودن عضو گروه است. این الگوریتم در زیرگروه چرخه‌ایِ متناظر با رمزگشایی بنالوح نیز قابل استفاده است؛ الزاماً به میدان اول محدود نیست، اما مرتبه و عملیات گروه باید درست مشخص شوند.

\section{کاربرد در بنالوح و تفاوت فضای پیام با پارامتر امنیت}
در بنالوح، پس از اعمال توان خصوصی، رابطۀ $a^m=b\pmod n$ در زیرگروهی با مرتبۀ $r$ به دست می‌آید. کوچک انتخاب شدن $r$ باعث می‌شود این لگاریتم گسسته برای رمزگشای مجاز قابل حل باشد. این مسئله با شکستن مستقیم متن رمز توسط کسی که عوامل $n$ را نمی‌داند یکسان نیست؛ توان خصوصی، متن رمز را به مسئلۀ محدودِ قابل حل تبدیل می‌کند.

در عین حال، بزرگ کردن $r$ برای افزایش ظرفیت مجموع وزن، هزینۀ همین حل را بالا می‌برد. بنابراین انتخاب $r$ باید هم قید $SW<r$ و هم حافظه و زمان رمزگشایی را رعایت کند. تغییر مقیاس اعشاری یا وزن‌ها نیز این مبادله را تغییر می‌دهد. این تحلیل، دلیل ضرورت بازآزمایی پارامترها هنگام اتصال بنچمارک کوچک به کلاینت‌های دارای وزن تعداد نمونه است.

الگوریتم شور روی رایانۀ کوانتومی مناسب، تهدیدی برای خانواده‌های متکی بر لگاریتم گسسته و تجزیه است. بزرگ کردن طول کلید این طرح‌های کلاسیک به‌تنهایی آن‌ها را پساکوانتومی نمی‌کند. زمان‌بندی عملی چنین تهدیدی به منابع و فناوری وابسته است و از بنچمارک کوچک این پایان‌نامه قابل استنتاج نیست.
